\begin{titlepage}
\thispagestyle{empty}
{\small УДК 004.056:343.721:159.9 (075.8)} \par
{\small \hspace*{2em} Є15} \par


\vspace{3mm}
{\footnotesize \itshape Погоджено з Українською міжбанківською асоціацією членів платіжних систем «ЄМА» (електронним листом від 06.01.2026 р.)} \par
%{\footnotesize За участі експертів Української міжбанківської асоціації членів платіжних систем «ЄМА».} \par

\vspace{3mm}
\textbf{Консультанти:} \par

%{\small Осадчук О. В., доктор технiчних наук, професор.} \par
%{\small Федоровська Р., керівниця EMA Academy та CyberSecurity School Асоціації «ЄМА»} \par
{\small \textbf{Данилов С.А.}, кандидат біологічних наук (PhD), нейробіолог, засновник науково-дослідного центру «Beehiveor Academy» та співзасновник освітньої платформи «Anima.ua». Фахівець у галузі нейрофізіології, психофізіології та поведінкової біології.} \par

%{\small \textbf{Вавілова А.}, кандидат педагогічних наук, клінічний психолог, нейропсихолог Психологічного центру розвитку людини «Зебра» (м. Київ). Фахівець із діагностики емоційної сфери, когнітивних функцій та психологічних наслідків стресу.} \par

{\small \textbf{Чубатюк І.І.}, магістр психології, психотерапевт (Вісбаденська академія позитивної психотерапії, WAPP, Німеччина), засновниця психологічної платформи «СЯЙВО», фахівець із сімейного консультування, кризової допомоги та особистісного розвитку.} \par

%{\small \textbf{Дзюба С.В.}, магістр філології, викладач англійської мови. Фахівець із професійної комунікації та перекладу спеціалізованої літератури.} \par

\vspace{3mm}
\textbf{Рецензенти:} \par
{\small \textbf{Петрова І.С.}, керівниця напряму розвитку бізнесу та корпоративних комунікацій Асоціації «ЄМА», директорка платформи Open API HUB. Фахівчиня з питань платіжної безпеки та протидії шахрайству.} \par

{\small \textbf{Гнатюк С.О.}, доктор технічних наук, професор, проректор з наукових досліджень та трансферу технологій КАІ, керівник науково-дослідної лабораторії протидії кіберзагрозам.} \par

{\small \textbf{Сліпухіна І.А.}, доктор педагогічних наук, професор, головний науковий співробітник відділу підтримки обдарованості Інституту обдарованої дитини НАПН України.} \par


\vspace{3mm}
{\footnotesize \textbf{\textit{Євдокимов С.О.}} Протидія шахрайству: Платформи оголошень та цифровi псевдопослуги.– Київ: Гуляєва В.М., 2026.– 300 с.} \par
\vspace{2mm}

{\footnotesize 

ISBN 978-617-8181-77-2 \par

DOI: https://doi.org/10.5281/zenodo.18314284


}

\vfill
{\small © Євдокимов С.О, 2026}
\end{titlepage}
\clearpage

\clearpage
\thispagestyle{empty} % Без номера сторінки
\null
\vfill % Центрування по вертикалі (початок)

\begin{center}
    % Обмежуємо ширину тексту до 75% від ширини сторінки для естетики
    \begin{minipage}{0.75\textwidth}
        \centering
        \itshape % Курсив
        \large   % Трохи більший шрифт
        \setlength{\parskip}{1.5em} % Відступ між "куплетами"
        
        % Встановлюємо трохи більший міжрядковий інтервал для легкості (опціонально)
        \linespread{1.2}\selectfont 

        На жаль, у сучасних реаліях, шахрайство перетворилося на високотехнологічну індустрію, 
        традиційні методи захисту втрачають свою ефективність.

        Найбільша помилка — думати, що шахраї грають за правилами.

        Інтернет-шахрай не сприймає людину як особистість.
        Його не цікавлять освіта, досвід чи соціальний статус.
        Для кіберзлочинця існує не людина, а можливість ---
        результат поведінки, помилок і необережних дій.

        Питання полягає лише в тому,
        чи створюють ці дії умови для перетворення користувача на жертву.
        
        \vspace{2em} % Відступ перед лінією
        
        % Елегантна горизонтальна лінія (довжина 30% від ширини тексту)
        % \hrule width 0.3\textwidth height 0.5pt 
    \end{minipage}
\end{center}

\vfill % Центрування по вертикалі (кінець)
\clearpage
